\section{第 11.1 节：对称性及其后果的概述}
\label{section:QM-C11-S01}
\studymeta
  {QM-C11-S01}
  {2026-10-12 至 2026-10-13}
  {Shankar, \textit{Principles of Quantum Mechanics}, Second Edition}
  {第 11 章 \S 11.1 Overview；印刷页 p.~279，PDF p.~289；至同页 \S 11.2 开始前}
  {讲解与概念核对已完成；笔记已确认}

\textbf{范围。}本节回顾经典力学中对称性的两个后果，并提出建立量子对应关系的目标。教材建议先复习\hyperref[section:QM-C02-S07]{\S 2.7 的正则变换}与\hyperref[section:QM-C02-S08]{\S 2.8 的对称性和守恒量}。本节没有编号 Example 或 Exercise，不另设题组；以下公式是对已学经典结论的简要展开，不提前推导 \S 11.2 的量子平移算符。

\notetopic{QM-C11-S01-T01}{同一动力学规律下的守恒量与运动解}
\keyidea{对称性要求变换保持系统的动力学规律，不要求每一个具体状态或每一条轨迹在变换后保持不变。}

\textbf{连续对称性与守恒量。}设 $g(q,p)$ 是不显含时间的 generator（生成元），它产生的 infinitesimal canonical transformation（无穷小正则变换）为
\begin{equation}
  \delta q_i=\epsilon\frac{\partial g}{\partial p_i},\qquad
  \delta p_i=-\epsilon\frac{\partial g}{\partial q_i}.
\end{equation}
这里 $\epsilon$ 是变换参数，$\delta$ 表示关于 $\epsilon$ 的一阶变化；$\{\ ,\ \}$ 表示 Poisson bracket（泊松括号）。在固定时刻，Hamiltonian 的变化为
\begin{align}
  \delta\mathcal H
  &=\sum_i\left(
      \frac{\partial\mathcal H}{\partial q_i}\delta q_i
     +\frac{\partial\mathcal H}{\partial p_i}\delta p_i\right)\notag\\
  &=\epsilon\sum_i\left(
      \frac{\partial\mathcal H}{\partial q_i}\frac{\partial g}{\partial p_i}
     -\frac{\partial\mathcal H}{\partial p_i}\frac{\partial g}{\partial q_i}\right)
   =\epsilon\{\mathcal H,g\}.
\end{align}
若该变换保持 $\mathcal H$ 不变，则 $\{\mathcal H,g\}=0$。沿真实运动轨迹，由 Hamilton 方程及括号的反对称性得到
\begin{equation}
  \boxed{\frac{\dd g}{\dd t}=\{g,\mathcal H\}
    =-\{\mathcal H,g\}=0}.
  \label{eq:QM-C11-S01-conserved-generator}
\end{equation}
即\textbf{连续对称性的生成元是守恒量}。若生成元显含时间，必须恢复 $\dd g/\dd t=\partial_tg+\{g,\mathcal H\}$，不能仅由括号为零断言守恒。离散对称性也不能直接套用这段无穷小生成元论证。

\textbf{对称变换把解映射成解。}设 $z=(q,p)$ 为相空间点，$T$ 是不显含时间的正则变换，且在每一时刻满足 $\mathcal H(Tz,t)=\mathcal H(z,t)$。正则性保持 Hamilton 方程的结构，Hamiltonian 不变则保证变换前后对应\textbf{同一个动力学系统}，因此
\begin{equation}
  \boxed{z(t)\text{ 是运动方程的解}
    \quad\Longrightarrow\quad Tz(t)\text{ 也是同一系统的解}}.
\end{equation}
两个条件缺一不可：只保持 Hamiltonian 的数值，并不足以保证运动方程被保持。新轨迹可以具有不同初始条件，并不必与原轨迹重合。

\textbf{对话补充：平移的直观含义。}一维自由粒子的 $\mathcal H=p^2/(2m)$ 不依赖位置。取 $g=p$，便有 $\delta q=\epsilon$、$\delta p=0$，所以平移不改变 Hamiltonian，动量守恒。原轨迹 $q(t)=q_0+p_0t/m$ 平移后成为 $q(t)+a$，仍满足同样的自由运动方程，但位置不同。在量子语言中，局域波包的整体平移同样可以改变物理状态；系统的平移对称性不等于波包必须平移不变。

\textbf{量子问题的入口。}本章随后要说明量子态的变换如何实现、如何判断 Hamiltonian 的对称性，以及它如何产生守恒量并将 Schr\"odinger 方程的解映射成另一解。\S 11.1 只提出这一目标，尚未建立具体的量子变换公式。

\textbf{一分钟回顾。}本节的两个经典结论是：不显含时间的连续对称性生成元守恒；保持 Hamiltonian 的不含时正则变换把运动解映射成同一系统的另一解。对称性描述规律不变，不要求状态不变；第 11 章将建立它们的量子对应关系。

\section{第 11.2 节：量子理论中的平移不变性}
\label{section:QM-C11-S02}
\studymeta
  {QM-C11-S02}
  {2026-10-13 至 2026-10-14}
  {Shankar, \textit{Principles of Quantum Mechanics}, Second Edition}
  {第 11 章 \S 11.2；印刷页 pp.~279--294，PDF pp.~289--304；至 \S 11.3 开始前}
  {五块讲解与习题 11.2.1、11.2.2 已核对；笔记已确认}

\textbf{主线。}先定义怎样平移量子态，再确定其生成元为动量，最后判断这种变换是否为系统的对称性。两项后果分别是动量守恒，以及平移后的初态按同样规律演化。教材关于经典正则变换的两段选读保留核心对应关系；11.2.3 未选做。完整习题和教材的原子--电容器示例见课程题目；不收录课堂概念检查。

\textbf{适用范围。}先讨论全实轴上无内部自由度的粒子，采用 $P=-\ii\hbar\partial_x$ 与 $[X,P]=\ii\hbar I$。无界算符的导数、乘积和对易子均在适当的共同定义域上使用；对称性还须保持定义域和边界条件。下文的指数由自伴生成元定义，不能仅凭微分表达式形式上的对易就忽略这些条件。

\notetopic{QM-C11-S02-T01}{平移量子态：主动与被动描述}
\textbf{来源：}印刷页 pp.~279--282，PDF pp.~289--292。
\keyidea{位置平移不是把粒子置于一个确定位置，而是将整个状态的位置分布刚性移动，保持动量分布不变。}

\textbf{从期望值到算符条件。}取 $a>0$ 表示向右平移。对任意归一化态，要求
\begin{equation}
  \ket{\psi_a}=T(a)\ket\psi,\qquad
  \expect{X}_a=\expect{X}+a,\qquad
  \expect{P}_a=\expect{P}.
\end{equation}
将变换移到算符上，得到等价条件
\begin{equation}
  \boxed{T^\dagger(a)XT(a)=X+aI,\qquad
         T^\dagger(a)PT(a)=P}.
  \label{eq:QM-C11-S02-transformed-XP}
\end{equation}
这里的“任意态”不可省略：一个 Hermitian 差算符若对所有态的期望都为零，其矩阵元也为零，才可推出算符本身为零；单个态的零期望不能这样使用。

主动描述保持测量坐标不动，把状态向右移 $a$；被动描述保持状态不动，把原点向左移 $a$，同一点的新坐标增加 $a$。两种描述通过
\begin{equation}
  \bra{\psi_a}\Omega\ket{\psi_a}
  =\bra\psi T^\dagger(a)\Omega T(a)\ket\psi
\end{equation}
给出相同的数值。注意这是一对\textbf{方向相反}的操作；原点向右移 $a$ 对应的坐标变化是 $X-aI$。

\textbf{位置表象中的作用。}先选取不附加位置相关相位的约定
\begin{equation}
  T(a)\ket x=\ket{x+a}.
\end{equation}
插入位置完备关系，有
\begin{align}
  \psi_a(x)
    &=\bra x T(a)\ket\psi
      =\int_{\real}\braket{x}{y+a}\psi(y)\dd y\notag\\
    &=\int_{\real}\delta(x-y-a)\psi(y)\dd y
      =\boxed{\psi(x-a)}.
\end{align}
若原波包中心是 $x_0$，新峰值满足 $x-a=x_0$，所以位于 $x_0+a$。自变量里的减号对应分布向右移动，并不矛盾。

平移保持整个内积：
\begin{equation}
  \int_{\real}\phi^*(x-a)\psi(x-a)\dd x
   =\int_{\real}\phi^*(y)\psi(y)\dd y.
\end{equation}
又因反向平移是其逆，$T$ 是 unitary（酉）算符。归一化和位置标准差不变，而位置期望增加 $a$。到这里构造的是一种变换，\textbf{尚未假定 Hamiltonian 在它作用下不变}。

\notetopic{QM-C11-S02-T02}{相位自由与动量生成元}
\textbf{来源：}印刷页 pp.~282--285，PDF pp.~292--295。

\textbf{为什么还要检查动量？}位置概率不能识别相位。只要求位置分布平移时，仍允许
\begin{equation}
  T_g(\epsilon)\ket x
   =\ee^{\ii\epsilon g(x)/\hbar}\ket{x+\epsilon},
  \qquad g(x)\in\real,
\end{equation}
对应
\begin{equation}
  \psi_\epsilon(x)
   =\ee^{\ii\epsilon g(x-\epsilon)/\hbar}\psi(x-\epsilon),
  \qquad |\psi_\epsilon(x)|^2=|\psi(x-\epsilon)|^2.
\end{equation}
然而动量包含空间导数，相位也被求导，因而
\begin{equation}
  \boxed{\expect P_\epsilon=\expect P
      +\epsilon\expect{g'(X)}}.
\end{equation}
完整推导见\relatedexercise{QM-C11-S02-E01}{教材习题 11.2.1}。要对\emph{所有}允许的初态保持平均动量，必须有 $g'(x)=0$，故 $g$ 仅为常数；它只产生整体相位，可约定为零。固定 $P=-\ii\hbar\partial_x$ 时，位置相关相位不能当作无关紧要的整体相位删掉。

\textbf{从小位移读出生成元。}写成
\begin{equation}
  T(\epsilon)=I-\frac{\ii\epsilon}{\hbar}G+O(\epsilon^2).
\end{equation}
其中 generator（生成元）$G$ 描述状态对平移参数的一阶响应。酉性的一阶条件给出 $G^\dagger=G$，计算见\relatedexercise{QM-C11-S02-E02}{教材习题 11.2.2}。在位置表象中比较
\begin{align}
  \bra x T(\epsilon)\ket\psi
    &=\psi(x)-\frac{\ii\epsilon}{\hbar}\bra xG\ket\psi
       +O(\epsilon^2),\\
  \psi(x-\epsilon)
    &=\psi(x)-\epsilon\psi'(x)+O(\epsilon^2),
\end{align}
便有
\begin{equation}
  \bra xG\ket\psi=-\ii\hbar\psi'(x)=\bra xP\ket\psi,
  \qquad \boxed{G=P}.
\end{equation}
这解释了动量算符与空间导数的联系：\textbf{动量生成空间平移}，不是因为平移会增加动量。

\textbf{从被动条件再推一次。}对任意算符 $\Omega$，一阶展开为
\begin{equation}
  T^\dagger\Omega T
   =\Omega+\frac{\epsilon}{\ii\hbar}[\Omega,G]
       +O(\epsilon^2).
\end{equation}
代入式~\eqref{eq:QM-C11-S02-transformed-XP}，得到
\begin{equation}
  [X,G]=\ii\hbar I,\qquad [P,G]=0.
\end{equation}
第一式与 $[X,P]=\ii\hbar I$ 比较，说明 $G-P$ 与 $X$ 对易；在这里的无内部自由度位置表示中，可写成 $G=P+F(X)$。再由
\begin{equation}
  0=[P,G]=[P,F(X)]=-\ii\hbar F'(X)
\end{equation}
知 $F$ 为常数。去掉对应的整体相位，仍得 $G=P$。只用第一条位置条件不能唯一确定生成元，还需要动量条件。

\notetopic{QM-C11-S02-T03}{平移不变性与动量守恒的层次}
\textbf{来源：}印刷页 pp.~283--289，PDF pp.~293--299；单态与全体态的区别为对话补充。
\keyidea{能定义平移算符，不等于系统具有平移对称性。对称性要求变换保持同一个 Hamiltonian。}

\textbf{系统的对称性条件。}要求对任意态均有
\begin{equation}
  \bra{\psi_a}H\ket{\psi_a}=\bra\psi H\ket\psi
  \quad\Longleftrightarrow\quad
  \boxed{T^\dagger(a)HT(a)=H}.
\end{equation}
酉共轭保持乘积次序。例如 $T^\dagger ABT=(T^\dagger AT)(T^\dagger BT)$，所以对已固定算符排序的 $H(X,P)$，
\begin{equation}
  T^\dagger H(X,P)T=H(X+aI,P).
\end{equation}
于是平移对称性就是 $H(X+aI,P)=H(X,P)$。对于通常的单粒子模型 $H=P^2/(2m)+V(X)$，任意距离的连续平移对称要求 $V(x+a)=V(x)$ 对所有 $a$ 成立，因此 $V$ 为常数。只对某个固定周期成立是\emph{离散}平移对称，不能据此推出连续平移的动量守恒结论。

取无穷小平移，
\begin{align}
  T^\dagger(\epsilon)HT(\epsilon)
   &=\left(I+\frac{\ii\epsilon}{\hbar}P\right)
        H\left(I-\frac{\ii\epsilon}{\hbar}P\right)
      +O(\epsilon^2)\notag\\
   &=H+\frac{\ii\epsilon}{\hbar}[P,H]+O(\epsilon^2).
\end{align}
与不变条件比较得到 $[P,H]=0$。由 Ehrenfest 定理，
\begin{equation}
  \boxed{\frac{\dd}{\dd t}\expect P
     =\frac1{\ii\hbar}\expect{[P,H]}=0}.
\end{equation}
对上述标量势模型，这也就是 $\dd\expect P/\dd t=-\expect{V'(X)}$。

\Needspace{4\baselineskip}
\textbf{必须分清的三个层次。}
\begin{enumerate}
  \item 对某一个态，$\expect{[P,H]}_{t_0}=0$ 只保证\emph{该时刻}平均动量的导数为零；要在一段时间内守恒，必须在整段演化上成立。
  \item 对某条演化轨迹，始终有 $\expect{[P,H]}_t=0$，即可保证该态的平均动量恒定，不要求算符 $[P,H]$ 为零。
  \item 系统具有连续平移对称性，则在前述适用条件下 $[P,H]=0$，并保证\emph{任意初态}的动量分布保持，不仅是某个态的一阶矩。
\end{enumerate}
例如谐振子满足 $[P,H]=-\ii\hbar m\omega^2X\ne0$，但基态始终有 $\expect X=\expect P=0$。这个单态的平均动量恒为零，并不能证明谐振子具有平移对称性。

\textbf{为何不止平均值守恒？}若 $H$ 不含时且与 $P$ 相容，则
\begin{equation}
  U(t)=\ee^{-\ii Ht/\hbar},\qquad [P,U(t)]=0.
\end{equation}
若初态是动量本征态，在广义本征态意义下
\begin{equation}
  PU(t)\ket p=U(t)P\ket p=p\,U(t)\ket p,
\end{equation}
所以本征值仍为同一个 $p$。对一般正规化态，$U$ 同样与动量的谱投影对易，任意动量区间的测量概率保持不变。含时 $H(t)$ 的情形，只要每个时刻都有相同的平移对称性，结论仍成立，但此时不能直接把演化写成 $\ee^{-\ii tH(t)/\hbar}$。

\notetopic{QM-C11-S02-T04}{有限平移与三种指数算符的区别}
\textbf{来源：}印刷页 pp.~287--291，PDF pp.~297--301；动量平移和时间演化的并列比较为对话补充。

\textbf{累积小平移。}将位移拆成许多小段，平移群的复合关系要求
\begin{equation}
  T(a+\epsilon)=T(\epsilon)T(a),\qquad T(0)=I.
\end{equation}
在生成元定义域上比较一阶项，得到
\begin{equation}
  \frac{\dd T(a)}{\dd a}=-\frac{\ii}{\hbar}PT(a),
  \qquad \boxed{T(a)=\ee^{-\ii aP/\hbar}}.
\end{equation}
因此
\begin{equation}
  T(a)T(b)=T(a+b),\qquad T^\dagger(a)=T(-a)=T^{-1}(a).
\end{equation}
位置表象中可形式地写作 $\ee^{-a\partial_x}$。当 Taylor 级数适用时，
\begin{equation}
  \ee^{-a\partial_x}\psi(x)
    =\psi(x)-a\psi'(x)+\frac{a^2}{2!}\psi''(x)+\cdots
    =\psi(x-a).
\end{equation}
\textbf{严格性补充。}自伴动量的谱定义使 $T(a)$ 在整个 $L^2(\real)$ 上有意义；平移公式不要求每个波函数都有收敛的 Taylor 级数。不能反过来把形式导数级数当成对任意态都成立的定义。

\textbf{在动量基中，平移只是逐分量加相位。}
\begin{equation}
  T(a)\ket p=\ee^{-\ii ap/\hbar}\ket p,
  \qquad \widetilde\psi_a(p)=\ee^{-\ii ap/\hbar}\widetilde\psi(p).
\end{equation}
这里指数里的 $p$ 是本征值，不再是算符。单个 $\ket p$ 仅获得整体相位；不同动量分量却获得不同相位，所以一般波包会改变。动量密度 $|\widetilde\psi(p)|^2$ 不变，当然也不是把 $\expect P$ 变为 $\expect P+p$。

\textbf{对话补充：位置平移、动量平移、时间演化。}令
\begin{equation}
  T(a)=\ee^{-\ii aP/\hbar},\qquad
  B(p_0)=\ee^{\ii p_0X/\hbar},\qquad
  U(t)=\ee^{-\ii tH/\hbar}\quad(H\text{ 不含时}).
\end{equation}
参数 $a$ 是位移，$p_0$ 是动量增量，$t$ 是经过的时间；它们不是把位置、动量或能量固定为参数值。
\begin{center}
\small
\renewcommand{\arraystretch}{1.3}
\begin{tabularx}{\textwidth}{@{}lXX@{}}
  \toprule
  变换 & 位置表象 & 动量表象 \\
  \midrule
  $T(a)$ & $\psi(x-a)$ & $\ee^{-\ii ap/\hbar}\widetilde\psi(p)$ \\
  $B(p_0)$ & $\ee^{\ii p_0x/\hbar}\psi(x)$ & $\widetilde\psi(p-p_0)$ \\
  $U(t)$ & 由具体 $H$ 决定 & 由具体 $H$ 决定 \\
  \bottomrule
\end{tabularx}
\end{center}
前两行的互换来自 Fourier 变换。例如
\begin{equation}
  \frac1{\sqrt{2\pi\hbar}}\int_{\real}
    \ee^{-\ii px/\hbar}\ee^{\ii p_0x/\hbar}\psi(x)\dd x
     =\widetilde\psi(p-p_0).
\end{equation}
相应地，
\begin{align}
  T^\dagger XT&=X+aI,& T^\dagger PT&=P,\\
  B^\dagger XB&=X,& B^\dagger PB&=P+p_0I.
\end{align}
所以 $B$ 不改位置密度，但使 $\expect P$ 增加 $p_0$；这不是测量，也不把动量分布压缩成一个本征值。两个指数的符号不同，保证所定义的位移和动量增量都取正方向，与 $[X,P]=\ii\hbar I$ 一致。

时间演化则在能量基中给每个分量乘上 $\ee^{-\ii Et/\hbar}$。自由粒子尤其能看清区别：
\begin{equation}
  \widetilde\psi(p,t)=\ee^{-\ii p^2t/(2m\hbar)}\widetilde\psi(p,0).
\end{equation}
位置平移的相位对 $p$ 是线性的，保持形状；自由演化的相位对 $p$ 是二次的，一般会导致波包展宽，不能用某个统一位移代替。只有特殊模型，如 $H=vP$，才有 $U(t)=T(vt)$。\textbf{指数形式相似，只说明都有生成元，不说明实施的是同一种变换。}

\textbf{教材选读：与正则变换的对应。}对 $U_G(\epsilon)=I-\ii\epsilon G/\hbar+O(\epsilon^2)$，
\begin{equation}
  \delta\Omega=\frac{\epsilon}{\ii\hbar}[\Omega,G]
  \quad\longleftrightarrow\quad
  \delta\omega=\epsilon\{\omega,g\}.
\end{equation}
酉共轭还保持对易关系，因为
\begin{equation}
  [U^\dagger X_iU,U^\dagger P_jU]
    =U^\dagger[X_i,P_j]U=\ii\hbar\delta_{ij}I.
\end{equation}
这对应经典正则变换保持基本 Poisson brackets；但保持正则结构还不等于保持 $H$，因此酉变换也不自动成为动力学对称性。这里是结构对应，不是任意经典函数都可无排序歧义地精确量子化的保证。

有限变换可展开为嵌套对易子：
\begin{equation}
  T^\dagger(a)\Omega T(a)
   =\Omega+\frac{a}{\ii\hbar}[\Omega,P]
     +\frac{a^2}{2!(\ii\hbar)^2}[[\Omega,P],P]+\cdots.
\end{equation}
经典对应式是 $\omega+a\{\omega,p\}+a^2\{\{\omega,p\},p\}/2!+\cdots$。各式按其级数适用条件理解；本次保留这一联系，不另做 11.2.3 的旋转级数练习。

\notetopic{QM-C11-S02-T05}{整体平移、总动量与相同的演化规律}
\textbf{来源：}印刷页 pp.~292--294，PDF pp.~302--304。

\textbf{多粒子的生成元。}整体平移要求所有位置同时增加相同的 $a$，所以
\begin{align}
  (T(a)\Psi)(x_1,\ldots,x_N)
    &=\Psi(x_1-a,\ldots,x_N-a),\\
  \Psi(x_1-\epsilon,\ldots,x_N-\epsilon)
    &=\Psi-\epsilon\sum_i\partial_{x_i}\Psi+O(\epsilon^2).
\end{align}
与无穷小生成元比较，得到
\begin{equation}
  \boxed{P_{\mathrm{tot}}=\sum_iP_i,\qquad
      T(a)=\ee^{-\ii aP_{\mathrm{tot}}/\hbar}}.
\end{equation}
有内部相互作用的系统仍可平移不变。例如
\begin{equation}
  H=\sum_i\frac{P_i^2}{2m_i}+\sum_{i<j}V_{ij}(X_i-X_j)
\end{equation}
在整体平移下，位置差变成 $(X_i+aI)-(X_j+aI)=X_i-X_j$，故 $T^\dagger HT=H$，总动量守恒。单个粒子的动量可以通过相互作用改变，不能推出每个 $P_i$ 都守恒。只移动一个粒子则通常改变位置差，不是这个系统的对称变换。

三维中相应地有 $T(\mathbf a)=\exp[-\ii\mathbf a\cdot\mathbf P_{\mathrm{tot}}/\hbar]$。沿哪个方向具有连续平移对称，就得到哪个总动量分量的守恒。

\textbf{对称性将动力学解映射成另一解。}若 $T^\dagger HT=H$，则 $[T,H]=0$；对不含时 $H$，进一步有 $[T,U(t)]=0$。因而
\begin{equation}
  \boxed{U(t)T(a)\ket{\psi(0)}
        =T(a)U(t)\ket{\psi(0)}}.
\end{equation}
两份初态只差一个位移，则之后始终只差同一个位移，而不是各自的波函数必须静止或不展宽。下图按教材图 11.1 的关系重新绘制：
\begin{center}
\begin{tikzpicture}[>=Stealth, every node/.style={font=\small}]
  \node (a) at (0,0) {$\ket{\psi(0)}$};
  \node (b) at (7,0) {$T(a)\ket{\psi(0)}$};
  \node (c) at (0,-1.6) {$U(t)\ket{\psi(0)}$};
  \node (d) at (7,-1.6) {$T(a)U(t)\ket{\psi(0)}$};
  \draw[->] (a) -- node[above]{平移 $T(a)$} (b);
  \draw[->] (a) -- node[left]{演化 $U(t)$} (c);
  \draw[->] (b) -- node[right]{演化 $U(t)$} (d);
  \draw[->] (c) -- node[below]{平移 $T(a)$} (d);
\end{tikzpicture}
\end{center}
这才是“同一个实验换个地点仍给出同样结果”的动力学内容；量子语境中的同样结果是\textbf{相应测量的概率分布相同}，不是两次随机测量必定取得同一个数。

\textbf{外部环境与空间均匀性。}若仅移动研究对象，却把对它施力的装置留在原处，对象与装置的相对位置已改变，一般不再是相同的实验。把相关场源也纳入系统并整体平移，才能讨论完整系统的空间均匀性。教材原子--电容器的具体模型见\relatedexercise{QM-C11-S02-E03}{教材正文示例}。不能简单断言“只要有外场就没有任何平移对称性”，应检查实际外势、平移方向和边界条件。

\textbf{一分钟回顾。}位置平移 $T(a)=\ee^{-\ii aP/\hbar}$ 由动量生成，但不增加动量；动量平移 $B(p_0)=\ee^{\ii p_0X/\hbar}$ 才增加 $p_0$。系统还需满足 $T^\dagger HT=H$ 才具有平移对称性，并由此得到动量守恒及 $TU=UT$。某个态的零对易子期望不能代替全系统的算符对易；多粒子整体平移对应的是总动量。
